\documentclass[12pt]{article}
\usepackage[margin=1in]{geometry}
\usepackage[T1]{fontenc}
\usepackage{array}
\usepackage{url}
\pagestyle{plain}

\newcommand{\transcriptionline}[2]{%
  \begin{center}
  \setlength{\tabcolsep}{0.45em}
  \renewcommand{\arraystretch}{1.15}
  \begin{tabular}{*{16}{>{\centering\arraybackslash}p{1.55em}}}
    #1 \\
    #2 \\
  \end{tabular}
  \end{center}
}

\begin{document}

\section*{Approximate Transcription}
\noindent Source recording: \url{signal-2026-09-15-08-11-29-091.aac}

This is an automatically derived, approximate transcription from the AAC recording. The melody is summarized below as note names with aligned lyric syllables, avoiding any external engraving dependency during compilation.

\medskip
\noindent Approximate setting: C major, common time, bass register.

\transcriptionline
  {C & C & D & E & G & G & E & D & C & D & E & D & C & & &}
  {Dad- & dy & and & Zo- & ey & are & go- & ing & for & a & morn- & ing & walk. & & &}

\transcriptionline
  {C & C & D & E & G & G & E & D & C & D & E & D & C & & &}
  {Dad- & dy & and & Zo- & ey & are & go- & ing & for & a & morn- & ing & walk. & & &}

\transcriptionline
  {E & E & G & G & A & A & G & F & E & E & D & D & C & C & D & C}
  {There & is & sun. & There & is & fun. & There's & some & bus- & i- & ness & that's & gon- & na & get & done.}

\transcriptionline
  {C & C & D & E & G & G & E & D & C & D & E & D & C & & &}
  {Dad- & dy & and & Zo- & ey & are & go- & ing & for & a & morn- & ing & walk. & & &}

\end{document}
